% formalisation.tex: blueprint extension of an accepted paper (AGENTS.md, section 6.6, case c).
% Not refereed: Lean checks every node. paper.tex is never edited.
% Nodes of paper.tex are referred to by their labels: xr-hyper reads paper.aux, so compile
% paper.tex first (CI does). New labels start with f (flem:, fprop:, fdef:) and are never reused.
% Every node has \lean{...} with its Lean name and an exhaustive \uses{...}; never write \leanok.
% Compile: latexmk -pdf -interaction=nonstopmode -halt-on-error -cd papers/<P>/formalisation.tex
\documentclass[11pt]{article}
\usepackage[T1]{fontenc}
\usepackage[utf8]{inputenc}
\usepackage{lmodern}
\usepackage[a4paper,margin=2.5cm]{geometry}
\usepackage{amsmath,amssymb,mathtools}
\usepackage{xcolor}
\usepackage{xr-hyper}
\usepackage[colorlinks=true,linkcolor=blue!50!black,citecolor=blue!50!black,urlcolor=blue!50!black]{hyperref}
\usepackage{../../templates/paper/blueprint-macros}
\externaldocument{paper}
\counterwithout{theorem}{section}
\renewcommand{\thetheorem}{F\arabic{theorem}}

% The paper's macros, copied from paper.tex (without the notation links).
\newcommand{\xmax}{x_{\max}}
\newcommand{\Simp}{\mathcal{P}_{S}}
\newcommand{\mean}{m}
\newcommand{\KL}{\mathrm{KL}}
\newcommand{\TV}{\mathrm{TV}}
\newcommand{\Ks}{\mathcal{K}_{S}}
\newcommand{\dirac}{\mathbf{1}_{\{x_{\max}\}}}
\newcommand{\Wone}{W_{1}}
\newcommand{\dF}{d_{\mathcal F}}
\newcommand{\om}{\omega}
\newcommand{\Gd}{G}

\title{Formalisation notes for P002-transfer-drift-between-laws}
\author{Incognita agents\thanks{Not refereed: every node is checked in Lean. Numbers F1, F2, \dots{}
  are this file's; other numbers refer to the paper.}}
\date{2026-11-20}

\begin{document}
\maketitle

Throughout, $S\subset[0,1]$ is finite, $\xmax=\max S$, $l\ge0$ is a drift budget (the paper's
$L_a$), $\om_d(q')=\sup\{\mean(q)-\mean(q'):q\in\Simp,\ d(q,q')\le l\}$ and
$\Gd^d=\inf\{\KL(r',q'):q'\in\Simp,\ \mean(q')+\om_d(q')\ge\mu\}$ (Proposition~\ref{prop:threshold}).

\section{Extra lemmas}

\begin{lemma}[The reach of a distance between the mean gap and the mixture bound]\label{flem:reach-min}
\lean{Incognita.TransferBandits.reachGain_eq_min,
  Incognita.TransferBandits.mix_single_mem_probVecs}
\uses{def:laws, def:ipm}
Let $q'\in\Simp$, and let $d=\dF$ satisfy $\mean(q)-\mean(q')\le d(q,q')$ for every $q\in\Simp$ and
$d((1-t)q'+t\dirac,q')\le t(\xmax-\mean(q'))$ for every $t\in[0,1]$. Then
$\om_d(q')=\min(l,\xmax-\mean(q'))$.
\end{lemma}

\begin{proof}
\uses{def:laws, def:ipm}
Every admissible $q$ has $\mean(q)-\mean(q')\le d(q,q')\le l$ and $\mean(q)\le\xmax$. Conversely,
with $g=\min(l,\xmax-\mean(q'))$ and $t=g/(\xmax-\mean(q'))\in[0,1]$ ($t=0$ if $\mean(q')=\xmax$),
the mixture $q_t=(1-t)q'+t\dirac$ is in $\Simp$, has $\mean(q_t)-\mean(q')=t(\xmax-\mean(q'))=g$ and
$d(q_t,q')\le g\le l$.
\end{proof}

\begin{lemma}[$\Wone$ satisfies the hypotheses of Lemma~\ref{flem:reach-min}]\label{flem:reach-w1}
\lean{Incognita.TransferBandits.abs_vecMean_sub_le_ipm_w1Class,
  Incognita.TransferBandits.ipm_w1Class_mix_single_le,
  Incognita.TransferBandits.reachGain_w1Class}
\uses{def:laws, def:ipm, flem:reach-min}
For $q,q'\in\Simp$ and $t\ge0$, $|\mean(q)-\mean(q')|\le\Wone(q,q')$ and
$\Wone((1-t)q'+t\dirac,q')\le t(\xmax-\mean(q'))$; hence $\om_{\Wone}(q')=\min(l,\xmax-\mean(q'))$.
\end{lemma}

\begin{proof}
\uses{def:laws, def:ipm, flem:reach-min}
The function $y\mapsto y$ is a $1$-Lipschitz test function with values in $[0,1]$. For a
$1$-Lipschitz $f$, $\sum_yf(y)\big((1-t)q'+t\dirac-q'\big)(y)=t\sum_yq'(y)(f(\xmax)-f(y))$, of
absolute value at most $t\sum_yq'(y)(\xmax-y)=t(\xmax-\mean(q'))$. Conclude with
Lemma~\ref{flem:reach-min}.
\end{proof}

\begin{lemma}[The threshold from the reach]\label{flem:threshold-ks}
\lean{Incognita.TransferBandits.thresholdKL_eq_vecKs}
\uses{def:laws, def:ipm}
If $\om_d(q')=\min(l,\xmax-\mean(q'))$ for every $q'\in\Simp$ and $\mu\le\xmax$, then
$\Gd^d=\Ks(r',\mu-l)$.
\end{lemma}

\begin{proof}
\uses{def:laws, def:ipm}
$\mean(q')+\min(l,\xmax-\mean(q'))=\min(\mean(q')+l,\xmax)\ge\mu$ if and only if
$\mean(q')\ge\mu-l$, as $\mu\le\xmax$: the two infima run over the same set.
\end{proof}

\begin{lemma}[The reach of total variation]\label{flem:reach-tv}
\lean{Incognita.TransferBandits.reachGain_tvClass}
\uses{def:laws, def:ipm}
For $q'\in\Simp$ and every real $l$, $\om_{\TV}(q')=\sup\{\sum_yw(y)(\xmax-y):0\le w\le q',\
\sum_yw(y)\le l\}$ (both suprema of the empty set, $0$, when $l<0$).
\end{lemma}

\begin{proof}
\uses{def:laws, def:ipm, lem:ipm}
By Lemma~\ref{lem:ipm}(c), $\TV(q,q')=\sum_y(q-q')_+(y)=\sum_y(q'-q)_+(y)$. Given an admissible $q$,
$w=(q'-q)_+$ satisfies $0\le w\le q'$, $\sum w=\TV(q,q')\le l$ and
$\mean(q)-\mean(q')=\sum_yy\big((q-q')_+-w\big)(y)\le\xmax\sum w-\sum_yyw(y)=\sum_yw(y)(\xmax-y)$.
Given $w$, $q=q'-w+(\sum w)\dirac$ is in $\Simp$, has $\TV(q,q')\le\sum w\le l$ and
$\mean(q)-\mean(q')=\sum_yw(y)(\xmax-y)$. Both sets are bounded above and their suprema are
nonnegative, so each supremum is at most the other.
\end{proof}

\begin{lemma}[$\Ks>0$ below the level]\label{flem:ks-pos}
\lean{Incognita.TransferBandits.ofReal_sq_sum_abs_sub_div_two_le_vecKL,
  Incognita.TransferBandits.ofReal_sq_sub_vecMean_div_two_le_vecKs,
  Incognita.TransferBandits.vecKs_pos, Incognita.TransferBandits.vecKs_lawVec_pos}
\uses{def:laws}
For $r,q\in\Simp$, $\KL(r,q)\ge\frac12\big(\sum_y|r(y)-q(y)|\big)^2$. Hence, if $\mean(r)<\ell$,
$\Ks(r,\ell)\ge(\ell-\mean(r))^2/2>0$.
\end{lemma}

\begin{proof}
\uses{def:laws}
The first inequality is Pinsker's inequality for probability vectors (when $\KL(r,q)<\infty$; it is
trivial otherwise). If $\mean(q)\ge\ell$, then
$\ell-\mean(r)\le\sum_yy(q(y)-r(y))\le\sum_y|q(y)-r(y)|$ as $S\subset[0,1]$.
\end{proof}

\begin{lemma}[A larger distance gives a smaller constant]\label{flem:order-mono}
\lean{Incognita.TransferBandits.mul_vecKs_le_programme,
  Incognita.TransferBandits.exists_one_le_programme,
  Incognita.TransferBandits.driftConstant_le_of_ipm_le}
\uses{def:ipm, def:constant}
For every pseudo-distance, $J_a(x,c;\mu)\ge x\Ks(\nu_a,\mu)$; if $\Ks(\nu_a,\mu)>0$, the set
$\{x>0:J_a(x,c;\mu)\ge1\}$ is nonempty. If $d_1\le d_2$ on $\Simp^2$ and $\Ks(\nu_a,\mu)>0$, then
$D^{d_2}_a(c)\le D^{d_1}_a(c)$.
\end{lemma}

\begin{proof}
\uses{def:ipm, def:constant}
The first law of a pair of $\Gamma_a(\mu)$ has mean at least $\mu$. The value $x=1/\Ks(\nu_a,\mu)$
(any $x>0$ if $\Ks=\infty$) is in the set. A pair feasible for $d_2$ is feasible for $d_1$, so
$J^{d_2}\ge J^{d_1}$ and the set defining $D^{d_1}_a(c)$ is contained in that of $D^{d_2}_a(c)$; it
is nonempty, and both are bounded below by $0$.
\end{proof}

\begin{lemma}[Abel summation bound]\label{flem:abel}
\lean{Incognita.TransferBandits.abs_sum_mul_sub_le,
  Incognita.TransferBandits.ipm_w1Class_le_ipm_ksClass,
  Incognita.TransferBandits.ipm_mono, Incognita.TransferBandits.ipm_ksClass_le_ipm_tvClass}
\uses{def:ipm}
Let $T\subset[m,b]$ be finite with $b\in T$, $f$ $1$-Lipschitz on $T$, and $D:T\to\mathbb R$ with
$|\sum_{y\in T,\,y\le s}D(y)|\le K$ for every $s\in\mathbb R$. Then
$|\sum_{y\in T}f(y)D(y)-f(b)\sum_{y\in T}D(y)|\le(b-m)K$. Consequently
$\Wone(q,q')\le\mathrm{KS}(q,q')\le\TV(q,q')$ for $q,q'\in\Simp$.
\end{lemma}

\begin{proof}
\uses{def:ipm}
Induction on the maximum: for $T'=T\cup\{b\}$ with $T<b$ and $b'=\max T$, the partial sums of $T$
are those of $T'$ at $\min(s,b')$, and
$\sum_{T'}fD-f(b)\sum_{T'}D=\big(\sum_TfD-f(b')\sum_TD\big)+(f(b')-f(b))\sum_TD$, bounded by
$(b'-m)K+(b-b')K$. For $\Wone$, apply it to $T=S$, $D=q-q'$ (so $\sum D=0$), $K=\mathrm{KS}(q,q')$,
$m=x_{\min}$, $b=\xmax$, and use $\xmax-x_{\min}\le1$. The indicators of $(-\infty,s]$ are test
functions of $\TV$.
\end{proof}

\begin{lemma}[Mixing the first law into the second]\label{flem:mix}
\lean{Incognita.TransferBandits.smul_add_smul_mem_probVecs,
  Incognita.TransferBandits.vecMean_smul_add_smul,
  Incognita.TransferBandits.vecKL_mix_le_mul, Incognita.TransferBandits.vecKL_mix_le,
  Incognita.TransferBandits.vecMean_mix_eq,
  Incognita.TransferBandits.exists_vecMean_eq_vecKL_le,
  Incognita.TransferBandits.ofReal_mul_vecKs_le}
\uses{def:laws}
For vectors $r,q$ and $t\in[0,1]$, $\KL(r,(1-t)q+tr)\le(1-t)\KL(r,q)$. Hence, for $r,q\in\Simp$ with
$\mean(r)\le\ell\le\mean(q)$, some $q_1\in\Simp$ has $\mean(q_1)=\ell$ and $\KL(r,q_1)\le\KL(r,q)$; and if
$\mean(r)<\mu\le\mean(q)$, then $(\mean(q)-\mean(r))\Ks(r,\mu)\le(\mu-\mean(r))\KL(r,q)$.
\end{lemma}

\begin{proof}
\uses{def:laws}
If $\KL(r,q)<\infty$, for $r(y)>0$ the concavity of $\log$ gives
$\log((1-t)q(y)+tr(y))\ge(1-t)\log q(y)+t\log r(y)$, so
$r(y)\log\frac{r(y)}{(1-t)q(y)+tr(y)}\le(1-t)\,r(y)\log\frac{r(y)}{q(y)}$; sum over $y$. If
$\KL(r,q)=\infty$ and $t<1$ the bound is trivial, and for $t=1$ both sides are $0$. For the
second claim, if $\mean(q)>\mean(r)$ take $q_1=(1-s)q+sr$ with
$s=(\mean(q)-\ell)/(\mean(q)-\mean(r))\in[0,1]$, which has mean $\ell$ (otherwise $q_1=q$). For the
third, the same mixture at level $\mu$ has $1-s=(\mu-\mean(r))/(\mean(q)-\mean(r))$ and
$\Ks(r,\mu)\le\KL(r,q_1)\le(1-s)\KL(r,q)$.
\end{proof}

\begin{lemma}[$\Ks<\infty$ below $\xmax$]\label{flem:ks-finite}
\lean{Incognita.TransferBandits.vecKL_ne_top, Incognita.TransferBandits.vecKs_ne_top}
\uses{def:laws}
If $q(y)>0$ wherever $r(y)>0$, then $\KL(r,q)<\infty$. For $r\in\Simp$ and $\ell<\xmax$,
$\Ks(r,\ell)<\infty$.
\end{lemma}

\begin{proof}
\uses{def:laws}
The first claim is the definition. If $\ell\le\mean(r)$, $\Ks(r,\ell)=0$. Otherwise the mixture
$(1-t)r+t\dirac$ with $t=(\ell-\mean(r))/(\xmax-\mean(r))\in(0,1)$ has mean $\ell$ and charges every
point that $r$ charges.
\end{proof}

\begin{lemma}[The programme of the mean gap]\label{flem:programme-mean}
\lean{Incognita.TransferBandits.nonempty_of_vecKs_ne_top,
  Incognita.TransferBandits.mul_biInf_of_nonempty,
  Incognita.TransferBandits.programme_meanClass,
  Incognita.TransferBandits.ofReal_driftConstant_meanClass_eq_ofReal,
  Incognita.TransferBandits.ofReal_driftConstant_meanClass}
\uses{def:laws, def:ipm, def:constant, flem:mix, flem:ks-pos}
Let $r,r'\in\Simp$, $l\ge0$, $\mean(r)\le\mu$, $\mean(r')\le\mu+l$, and $\Ks(r,\mu)$, $\Ks(r',\mu-l)$
finite. For the mean gap, $J(x,c;\mu)=x\Ks(r,\mu)+c\Ks(r',\mu-l)$ for all $x,c$. If moreover
$S\subset[0,1]$, $\mean(r)<\mu$ and $c\ge0$, then $D(c)=(1-c\Ks(r',\mu-l))_+/\Ks(r,\mu)$.
\end{lemma}

\begin{proof}
\uses{def:laws, def:ipm, def:constant, flem:mix, flem:ks-pos}
A pair of the mean gap has $\mean(q)\ge\mu$ and $\mean(q')\ge\mean(q)-l\ge\mu-l$, which gives $\ge$.
Conversely, let $\mean(q)\ge\mu$ and $\mean(q')\ge\mu-l$. By Lemma~\ref{flem:mix} replace $q$ by a law
of mean $\mu$, and, if $\mean(q')>\mu+l$, $q'$ by a law of mean $\mu+l$ (as $\mean(r')\le\mu+l$),
without increasing the divergences; the new pair has $|\mean(q_1)-\mean(q'_1)|\le l$. Taking the
infima over $q$ and $q'$ gives $\le$ (the two sets are nonempty since the $\Ks$ are finite). Then
$J(x,c)\ge1$ if and only if $x\ge(1-c\Ks(r',\mu-l))/\Ks(r,\mu)$, as $\Ks(r,\mu)>0$
(Lemma~\ref{flem:ks-pos}), and $D(c)$ is the positive part of this bound.
\end{proof}

\begin{lemma}[A source mean above the band]\label{flem:mean-high}
\lean{Incognita.TransferBandits.ofReal_sq_div_two_le_vecKL,
  Incognita.TransferBandits.ofReal_driftConstant_meanClass_lt}
\uses{def:laws, def:ipm, def:constant, flem:mix, flem:ks-pos}
Let $S\subset[0,1]$, $r,r'\in\Simp$, $\mean(r)<\mu$, $\Ks(r,\mu)<\infty$, $\mean(r')>\mu+l$ and $c>0$.
Then, for the mean gap, $D(c)<1/\Ks(r,\mu)$.
\end{lemma}

\begin{proof}
\uses{def:laws, def:ipm, def:constant, flem:mix, flem:ks-pos}
Let $k=\Ks(r,\mu)$, $g=\mu-\mean(r)$, $\delta=(\mean(r')-\mu-l)/2$ and $\eta=\delta^2/2$. Pinsker's
inequality (Lemma~\ref{flem:ks-pos}) gives $\KL(r',q')\ge\eta$ whenever
$\mean(q')\le\mean(r')-\delta$. Let $x=\max\big(1/(k(1+\delta/g)),(1-c\eta)/k\big)<1/k$. A pair of the
mean gap with $\mean(q)\ge\mu+\delta$ has $\KL(r,q)\ge k(1+\delta/g)$ by Lemma~\ref{flem:mix}, so its
cost is at least $1$; a pair with $\mean(q)<\mu+\delta$ has $\mean(q')\le\mean(q)+l<\mean(r')-\delta$,
so its cost is at least $xk+c\eta\ge1$. Hence $J(x,c)\ge1$ and $D(c)\le x<1/k$.
\end{proof}

\begin{lemma}[When the source helps the mean gap]\label{flem:criteria-mean}
\lean{Incognita.TransferBandits.ofReal_driftConstant_meanClass_lt_iff}
\uses{def:laws, def:ipm, def:constant, flem:programme-mean, flem:mean-high, flem:ks-pos}
Let $S\subset[0,1]$, $r,r'\in\Simp$, $l\ge0$, $\mean(r)<\mu$, $\Ks(r,\mu)$ and $\Ks(r',\mu-l)$ finite,
and $c>0$. For the mean gap, $D(c)<1/\Ks(r,\mu)$ if and only if $l<|\mean(r')-\mu|$.
\end{lemma}

\begin{proof}
\uses{def:laws, def:ipm, def:constant, flem:programme-mean, flem:mean-high, flem:ks-pos}
If $\mean(r')>\mu+l$, both sides hold (Lemma~\ref{flem:mean-high}). Otherwise, by
Lemma~\ref{flem:programme-mean}, $D(c)<1/\Ks(r,\mu)$ if and only if $c\Ks(r',\mu-l)>0$, that is
$\Ks(r',\mu-l)>0$, that is $\mean(r')<\mu-l$ (Lemma~\ref{flem:ks-pos}, and $\Ks(r',\mu-l)=0$ when
$\mean(r')\ge\mu-l$), which, as $\mean(r')\le\mu+l$, is $l<|\mean(r')-\mu|$.
\end{proof}

\begin{lemma}[A binding projection sits on its level]\label{flem:proj-level}
\lean{Incognita.TransferBandits.vecMean_eq_of_isKLProj}
\uses{def:laws, def:constant, flem:mix}
Let $\rho\in\Simp$ and let $Q$ minimise $\KL(\rho,\cdot)$ on $\{q\in\Simp:\mean(q)\ge\ell\}$. If
$0<\Ks(\rho,\ell)<\infty$, then $\mean(Q)=\ell$.
\end{lemma}

\begin{proof}
\uses{def:laws, def:constant, flem:mix}
As $\Ks(\rho,\ell)>0$, $\mean(\rho)<\ell$. If $\mean(Q)>\ell$, the mixture $(1-s)Q+s\rho$ of mean $\ell$
has $s>0$ and, by Lemma~\ref{flem:mix}, divergence at most $(1-s)\Ks(\rho,\ell)<\Ks(\rho,\ell)$, a
contradiction.
\end{proof}

\begin{lemma}[Total variation and the width of the support]\label{flem:tv-width}
\lean{Incognita.TransferBandits.mul_ipm_tvClass_sub_vecMean_sub,
  Incognita.TransferBandits.vecMean_sub_le_mul_ipm_tvClass,
  Incognita.TransferBandits.eq_of_vecMean_sub_eq_mul_ipm_tvClass}
\uses{def:laws, def:ipm, lem:ipm}
Let $S\subset[b_0,b_1]$ and $q,q'\in\Simp$. Then
$(b_1-b_0)\TV(q,q')-(\mean(q)-\mean(q'))=\sum_y\big((b_1-y)(q-q')_+(y)+(y-b_0)(q'-q)_+(y)\big)\ge0$; in
case of equality, $q(y)=q'(y)$ for every $y\in S\cap(b_0,b_1)$.
\end{lemma}

\begin{proof}
\uses{def:laws, def:ipm, lem:ipm}
By Lemma~\ref{lem:ipm}(c), $\TV(q,q')=\sum_y(q-q')_+(y)=\sum_y(q'-q)_+(y)$, and
$q-q'=(q-q')_+-(q'-q)_+$; expand. In case of equality every term vanishes, and for $y\in(b_0,b_1)$
both factors $b_1-y$ and $y-b_0$ are positive.
\end{proof}

\section{New decompositions of the paper's proofs}

\begin{proof}[Formal route to Proposition~\ref{prop:order}]
\proves{prop:order}
\uses{flem:ks-pos, flem:ks-finite, flem:programme-mean, flem:order-mono, flem:abel, lem:ipm, lem:ipm-order}
Part~(a) is Lemma~\ref{flem:programme-mean} with $r=\nu_a$, $r'=\nu'_a$, $\mu=\mu^\star$, $l=L_a$:
both $\Ks$ are finite by Lemma~\ref{flem:ks-finite}, as $\mu^\star-L_a<\mu^\star<\xmax$. The paper's
pair $(Q^\star,Q'^\star)$ is replaced by near-minimisers moved to the band
$|\mean(q)-\mean(q')|\le L_a$, so neither the existence nor the uniqueness of the projections is
needed. Parts~(b) and~(c) are as in the paper, with Lemma~\ref{flem:order-mono} for the step
``$d_1\le d_2$ gives $D^{d_2}\le D^{d_1}$'', which needs the set defining $D^{d_1}$ to be nonempty:
this is where $\Ks(\nu_a,\mu^\star)>0$ (Lemma~\ref{flem:ks-pos}) enters. The order
$|\mean(q)-\mean(q')|\le\Wone\le\mathrm{KS}\le\TV$ is Lemma~\ref{flem:abel} and the test function
$y\mapsto y$.
\end{proof}

\begin{proof}[Formal route to Proposition~\ref{prop:criteria}~(b) for the mean gap]
\proves{prop:criteria}
\uses{flem:ks-finite, flem:criteria-mean}
Lemma~\ref{flem:criteria-mean} with $r=\nu_a$, $r'=\nu'_a$, $\mu=\mu^\star$, $l=L_a$ (finite $\Ks$ by
Lemma~\ref{flem:ks-finite}); when $\mean(\nu'_a)\le\mu^\star+L_a$, $L_a<|\mean(\nu'_a)-\mu^\star|$ reads
$\mean(\nu'_a)<\mu^\star-L_a$. The paper's route (Lemma~\ref{lem:criterion-point} at $x_0=1/\Ks$, with
$d(Q^\star,\nu'_a)=|\mu^\star-\mean(\nu'_a)|$) is replaced by Lemmas~\ref{flem:programme-mean}
and~\ref{flem:mean-high}. Part~(a) and the first sentence of~(b) are not covered.
\end{proof}

\begin{proof}[Formal route to Proposition~\ref{prop:tv}~(a)]
\proves{prop:tv}
\uses{flem:ks-finite, flem:proj-level, flem:tv-width}
As $\mathcal K'_0=\Ks(\nu'_a,\mu^\star-L_a)>0$ and $\mathcal K'_0<\infty$
(Lemma~\ref{flem:ks-finite}), $\mean(Q'^\star)=\mu^\star-L_a$
(Lemma~\ref{flem:proj-level}), so $\mean(Q^\star)-\mean(Q'^\star)\ge L_a$. Lemma~\ref{flem:tv-width} with
$b_0=x_{\min}<b_1=\xmax$ (as $|S|\ge2$) gives $L_a\le(\xmax-x_{\min})\TV(Q^\star,Q'^\star)$, and in case
of equality in the bound of~(a), equality in Lemma~\ref{flem:tv-width}. Parts~(b) and~(c) are not
covered.
\end{proof}

\begin{proof}[Formal route to Proposition~\ref{prop:threshold}~(b) and~(c)]
\proves{prop:threshold}
\uses{flem:reach-min, flem:reach-w1, flem:threshold-ks, flem:reach-tv}
Part~(b) for the mean gap is \texttt{reachGain\_meanClass} (already in the library) and for $\Wone$
Lemma~\ref{flem:reach-w1}; the identity $\Gd^d=\Ks(\nu'_a,\mu^\star-L_a)$ follows for both by
Lemma~\ref{flem:threshold-ks}, as $\mu^\star<\xmax$. Part~(c) is Lemma~\ref{flem:reach-tv}. Part~(a)
is not covered by this route.
\end{proof}

\end{document}
